% Experiments and Results (full version: all tables generated by src/make_tables.py; hand-written numbers are forbidden)
\label{sec:experiments}

\noindent\textbf{Confirmatory versus exploratory analysis.} The test
set is used for one confirmatory evaluation---the head-to-head
comparison of Table~\ref{tab:main}---and for the financial protocol
described in Section~\ref{sec:protocol}. All other analyses
(ablations, permutation importance, difficulty stratification,
policy comparisons, error analysis, and LLM diagnostics) are
\emph{exploratory}: they generate hypotheses about the structure of
the task rather than testing pre-registered claims, and they should
be read as descriptive evidence on a single season rather than as
confirmatory tests. No threshold, weight, or hyperparameter is tuned
on the test set; all such choices are fixed on the validation split
before the test set is examined.

% ================= 4.1 overall performance =================
\subsection{Overall Performance}
\begin{table*}[tbp]
\centering
\small
\caption{Overall performance on the held-out test period (1,104 matches; 2025/26 season to date). The LLM parsed 100\% of matches at a cost of \$0.55. Bootstrap 95\% confidence intervals are reported in the text.}
\label{tab:main}
\resizebox{\textwidth}{!}{
\begin{tabular}{lccccccc}
\toprule
Model & Acc & LogLoss & Brier & ECE & ROI & MDD & Sharpe \\
\midrule
Market (de-vig) & 54.2\% & 0.967 & 0.575 & 0.026 & -0.4\% & -3.4\% & -0.004 \\
Dixon--Coles & 51.1\% & 0.996 & 0.594 & 0.023 & -2.0\% & -4.3\% & -0.019 \\
Poisson & 50.8\% & 1.003 & 0.596 & 0.045 & -2.3\% & -3.6\% & -0.022 \\
Elo & 51.6\% & 0.998 & 0.595 & 0.031 & -1.6\% & -3.9\% & -0.015 \\
Random Forest & 54.7\% & 0.973 & 0.577 & 0.042 & 2.8\% & -2.1\% & 0.028 \\
XGBoost & 54.2\% & 0.971 & 0.576 & 0.022 & 0.5\% & -2.7\% & 0.005 \\
LLM (DeepSeek) & 53.9\% & 0.967 & 0.575 & 0.025 & -1.0\% & -3.2\% & -0.011 \\
\bottomrule
\end{tabular}
}
\end{table*}


Table~\ref{tab:main} reports the overall performance on the
temporally held-out test period (2025/26 season to date). Several observations are immediate. First, the de-vigged
closing odds achieve 54.2\% accuracy and a log loss of 0.967, and no
learned model exceeds them in accuracy within statistical uncertainty
(bootstrap 95\% confidence intervals for accuracy all overlap the
market point estimate; the largest gap, Random Forest at 54.7\%, is
well inside the market's CI of [51.4\%, 57.2\%]). Second, the
domain-enhanced LLM reaches a log loss of 0.967 and ECE of 0.025,
statistically indistinguishable from the market (0.967 and 0.026):
a zero-training prompt reaches market-level probability quality.
The reliability diagram in Figure~\ref{fig:reliability} confirms
this: all three probability sources track the diagonal, with the LLM
and the market showing the smallest deviation.
Third, all ROI values are small and their confidence intervals
cross zero: no method, including the market itself, shows a stable
profit over this test period under equal-stake betting. This is the
expected regime of an efficient market, and it is the backdrop
against which all subsequent analyses are interpreted. The reader is
reminded that the held-out test period (2025/26 season to date) is
incomplete (matches through 2026-02-12): accuracy estimates are stable (they depend on many
games), but the financial point estimates above may shift when the
season completes, and the bootstrap intervals should be read
accordingly.

% ================= 4.2 ablation =================
\subsection{Ablation Study}
\begin{table}[tbp]
\centering
\small
\setlength{\tabcolsep}{5pt}
\caption{Feature-group ablation on the complete 2025/26 season. ROI is descriptive and uses the common closing-line settlement protocol.}
\label{tab:ablation}
\begin{tabular}{lcccc}
\toprule
Configuration & Features & Accuracy & Log loss & ROI \\
\midrule
Full (60 features) & 60 & 52.9\% & 0.983 & -1.9\% \\
Without market odds & 49 & 51.4\% & 1.004 & -2.3\% \\
Without xG & 53 & 52.5\% & 0.983 & -3.3\% \\
Without referee & 56 & 52.8\% & 0.983 & -2.5\% \\
Without rank & 58 & 52.3\% & 0.985 & -3.8\% \\
Without form & 30 & 52.7\% & 0.983 & -2.6\% \\
Structured non-market only & 45 & 51.1\% & 1.006 & -2.9\% \\
\bottomrule
\end{tabular}
\end{table}


Table~\ref{tab:ablation} ablates feature groups from XGBoost. Removing
market odds costs 3.1 percentage points of accuracy (54.2\% to
51.1\%) and 0.028 of log loss, confirming that bookmaker information
dominates. Removing referee statistics changes nothing (54.2\% vs
54.2\%): the referee features provide no measurable predictive gain,
a negative result reported deliberately. (One caveat: referee
records exist only for the Premier League in the present data, so the referee
feature group is evaluated on a subset of matches; the null result
should be read as \emph{no gain where referee data exist}, not as a
global statement about referee information.) xG features contribute a
small but consistent improvement (removing them raises log loss
slightly), and rank and recent-form features contribute little beyond
what the market already encodes.

% ================= 4.3 feature importance =================
\subsection{Feature Importance}
\begin{table}[htbp]
\centering
\small
\caption{Permutation importance by feature group (mean increase in log loss over 5 permutations, test set).}
\label{tab:importance}
\resizebox{\linewidth}{!}{
\begin{tabular}{lcc}
\toprule
Group & Importance & Share \\
\midrule
Market & 0.0745 & 88.3\% \\
xG & 0.0061 & 7.3\% \\
Season cumulative & 0.0027 & 3.2\% \\
League rank & 0.0005 & 0.6\% \\
Referee & 0.0005 & 0.6\% \\
Rolling form & 0.0001 & 0.1\% \\
\bottomrule
\end{tabular}
}
\end{table}


Permutation importance (Table~\ref{tab:importance}) quantifies the
same hierarchy: market-derived features account for roughly 87\% of
the aggregate permutation importance under the XGBoost specification,
xG for about 6\%, and referee,
rank, and form features for under 2\% combined. The result is robust
across the ablation and permutation analyses and motivates the
position of this paper: the marginal information left to any
forecasting model is small and concentrated in a few market-driven
signals.

% ================= 4.4 by league =================
\subsection{League and Season Generalization}
\begin{table}[htbp]
\centering
\small
\caption{Per-league performance of the global XGBoost model on the test set.}
\label{tab:byleague}
\resizebox{\linewidth}{!}{
\begin{tabular}{lccccc}
\toprule
League & N & Acc & LogLoss & ROI & MDD \\
\midrule
D1 & 188 & 57.4\% & 0.945 & 6.9\% & -6.0\% \\
E0 & 260 & 50.0\% & 1.016 & -7.4\% & -8.1\% \\
F1 & 189 & 55.0\% & 0.953 & 1.1\% & -6.8\% \\
I1 & 239 & 52.7\% & 0.981 & -2.4\% & -7.3\% \\
SP1 & 228 & 57.0\% & 0.947 & 7.0\% & -2.9\% \\
\bottomrule
\end{tabular}
}
\end{table}

\begin{table}[htbp]
\centering
\small
\caption{Leave-one-league-out generalization: model trained on the other four leagues, evaluated on the held-out league (test set).}
\label{tab:lolo}
\resizebox{\linewidth}{!}{
\begin{tabular}{lcccc}
\toprule
Held-out league & N & Acc & LogLoss & ROI \\
\midrule
E0 & 260 & 47.3\% & 1.015 & -13.5\% \\
D1 & 188 & 54.3\% & 0.942 & -0.9\% \\
F1 & 189 & 55.6\% & 0.958 & 4.5\% \\
I1 & 239 & 53.6\% & 0.979 & 0.6\% \\
SP1 & 228 & 53.5\% & 0.954 & -2.2\% \\
\bottomrule
\end{tabular}
}
\end{table}


Performance varies markedly across leagues (Table~\ref{tab:byleague}):
the Bundesliga and La Liga are the most predictable (57.4\% and
57.0\% accuracy), while the Premier League is the least (50.0\%), the
most liquid and informationally efficient of the five. The
leave-one-league-out experiment (Table~\ref{tab:lolo}) shows the same
ordering in a stricter setting: a model trained on the other four
leagues degrades most on the Premier League (47.3\%). These results
indicate that market efficiency, not model capacity, is the binding
constraint on accuracy.

% ================= 4.5 walk-forward =================
\subsection{Temporal Robustness}
\begin{table}[htbp]
\centering
\small
\caption{Walk-forward evaluation across four test seasons. Market: de-vigged closing odds. XGB-exp: expanding training window. XGB-roll2: last two seasons only.}
\label{tab:walkforward}
\resizebox{\linewidth}{!}{
\begin{tabular}{lccccc}
\toprule
Season & N & Market & XGB-exp & LL(exp) & XGB-roll2 \\
\midrule
2022/2023 & 1826 & 54.1\% & 52.5\% & 0.991 & 53.0\% \\
2023/2024 & 1752 & 55.9\% & 53.9\% & 0.968 & 54.0\% \\
2024/2025 & 1752 & 54.9\% & 54.1\% & 0.974 & 51.3\% \\
2025/2026 & 1104 & 54.2\% & 54.2\% & 0.971 & 53.5\% \\
\bottomrule
\end{tabular}
}
\end{table}


Table~\ref{tab:walkforward} evaluates the model in a walk-forward
manner over four test seasons. The market baseline is at least as
accurate as XGBoost in every season, and restricting training to the
last two seasons (XGB-roll2) does not help: there is no evidence of
``model aging'' within this seven-year window, and retraining
frequency has no measurable effect on accuracy. This addresses the
common concern that static models deteriorate as team dynamics
evolve; on these data, the information content of older seasons is not
harmful.

% ================= 4.6 draws =================
\subsection{Draw Prediction}
\begin{table}[htbp]
\centering
\small
\caption{Class-wise metrics (XGBoost, test set). Empirical draw rate is 25.5\% but the model predicts a draw in only 2.3\% of matches.}
\label{tab:drawclass}
\resizebox{\linewidth}{!}{
\begin{tabular}{lcccc}
\toprule
Class & Prec & Recall & F1 & ECE \\
\midrule
Home win & 0.566 & 0.818 & 0.669 & 0.037 \\
Draw & 0.360 & 0.032 & 0.059 & 0.007 \\
Away win & 0.508 & 0.567 & 0.536 & 0.020 \\
Pred. draw rate & 2.3\% & --- & --- & --- \\
Empirical draw rate & 25.5\% & --- & --- & --- \\
\bottomrule
\end{tabular}
}
\end{table}

\begin{table}[htbp]
\centering
\small
\caption{Cost-sensitive training: draw class weight sweep (XGBoost, test set). Weight 1.0 is the baseline.}
\label{tab:drawweight}
\resizebox{\linewidth}{!}{
\begin{tabular}{lccccc}
\toprule
w(draw) & Acc & Macro-F1 & Draw R & Draw P & LogLoss \\
\midrule
1.0 & 54.2\% & 0.421 & 0.032 & 0.360 & 0.971 \\
1.5 & 52.0\% & 0.493 & 0.356 & 0.317 & 0.987 \\
2.0 & 46.6\% & 0.456 & 0.584 & 0.300 & 1.012 \\
2.5 & 45.2\% & 0.443 & 0.722 & 0.303 & 1.037 \\
3.0 & 45.6\% & 0.447 & 0.612 & 0.301 & 1.047 \\
\bottomrule
\end{tabular}
}
\end{table}

\begin{table}[htbp]
\centering
\small
\caption{Draw diagnostics (test set). First five rows: bins of the model's predicted draw probability with the empirical draw rate in each bin. The model's draw probabilities are calibrated (class-level ECE 0.007) and monotone, but weakly discriminative (r = 0.12 vs the market's 0.14) and never the argmax at typical values near the base rate, which explains the near-empty draw column in \ref{fig:confusion}.}
\label{tab:drawdeep}
\resizebox{\linewidth}{!}{
\begin{tabular}{lccc}
\toprule
Model p(draw) bin & N & Mean p & Empirical draw rate \\
\midrule
(0,0.1] & 10 & 0.084 & 30.0\% \\
(0.1,0.2] & 225 & 0.161 & 16.9\% \\
(0.2,0.25] & 207 & 0.229 & 20.8\% \\
(0.25,0.3] & 407 & 0.276 & 28.0\% \\
(0.3,0.4] & 253 & 0.326 & 32.4\% \\
Class ECE: H / D / A & --- & 0.037 / 0.007 / 0.020 & --- \\
Class log loss: H / D / A & --- & 0.733 / 1.350 / 1.001 & --- \\
corr(p(draw), draw) & --- & 0.124 & --- \\
corr(market p(draw), draw) & --- & 0.137 & --- \\
\bottomrule
\end{tabular}
}
\end{table}


Draws are systematically under-predicted: the empirical draw rate is
25.5\%, yet the model predicts a draw in only 2.3\% of matches, and
draw recall is 0.032 (Table~\ref{tab:drawclass}). The confusion
matrix in Figure~\ref{fig:confusion} shows the near-empty draw
column visually. The market behaves
similarly (0.3\% predicted draw rate). A deeper analysis
(Table~\ref{tab:drawdeep}) shows that this is not a probability-
estimation failure: the model's draw probabilities are well
calibrated (class-level ECE of 0.007 for draws, versus 0.037 and
0.020 for home and away wins) and increase monotonically with the
empirical draw rate across the probability range. The problem is
discriminative, not probabilistic: draw probabilities correlate with
outcomes only weakly (r = 0.12, comparable to the market's 0.14),
and the draw class contributes the largest share of multiclass log
loss (1.35 vs 0.73 for home wins). Because draw probabilities
concentrate near 0.25---the empirical base rate---they are rarely the
argmax in a three-way competition, so recall collapses at the
decision level even though the underlying probability is sound. This
is a structural property of the task rather than a defect of any
single model. The per-class Brier decomposition (Table~\ref{tab:drawdeep})
localizes the failure precisely: for the draw class the uncertainty
term is fixed by the base rate (0.190), the model's resolution is
0.003---an order of magnitude below the 0.040 achieved for home
wins---and its reliability error is negligible (0.001); the market's
resolution (0.002) is only marginally worse. The draw probability is
calibrated and monotone, but it barely separates draws from
non-draws, and that is a ceiling the model inherits from the task.
Table~\ref{tab:drawweight} shows that cost-sensitive
training recovers balance in a controllable way: a draw weight of
1.5 raises macro-F1 from 0.421 to 0.493 and draw recall from 0.032
to 0.356, at the cost of 2.2 points of accuracy. The accuracy--balance
trade-off is monotone and easy to report.

% ================= 4.7 value betting =================
\subsection{Value Betting and Market Efficiency}
\begin{table}[htbp]
\centering
\small
\caption{Value betting: bet when EV = $p \times odds - 1$ exceeds the threshold (XGBoost probabilities, closing odds). Higher thresholds select increasingly confident bets and lose more money.}
\label{tab:value}
\resizebox{\linewidth}{!}{
\begin{tabular}{lccccc}
\toprule
EV thresh. & ROI & Sharpe & MDD & N bets & Coverage \\
\midrule
0.00 & -11.5\% & -0.077 & -14.5\% & 889 & 80.5\% \\
0.02 & -12.7\% & -0.084 & -15.9\% & 756 & 68.5\% \\
0.05 & -15.2\% & -0.097 & -19.1\% & 584 & 52.9\% \\
0.08 & -15.8\% & -0.095 & -19.9\% & 420 & 38.0\% \\
0.10 & -19.0\% & -0.113 & -25.7\% & 351 & 31.8\% \\
0.15 & -26.7\% & -0.145 & -32.8\% & 184 & 16.7\% \\
\bottomrule
\end{tabular}
}
\end{table}

\begin{table}[htbp]
\centering
\small
\caption{Value betting after isotonic calibration (fit on validation). Losses shrink but remain negative: calibration cannot create information the model does not have.}
\label{tab:valuecal}
\resizebox{\linewidth}{!}{
\begin{tabular}{lccc}
\toprule
EV thresh. & ROI & N bets & Coverage \\
\midrule
0.00 & -7.1\% & 965 & 87.4\% \\
0.02 & -5.3\% & 876 & 79.3\% \\
0.05 & -7.0\% & 724 & 65.6\% \\
0.08 & -7.7\% & 574 & 52.0\% \\
0.10 & -7.5\% & 484 & 43.8\% \\
\bottomrule
\end{tabular}
}
\end{table}


Table~\ref{tab:value} reports the value-betting experiment: bets are
placed when the model's expected value $p \times odds - 1$ exceeds a
threshold, using closing odds. Every threshold evaluated loses
money in the test period, and
losses grow with the threshold: the most ``confident'' bets are the
worst. After isotonic calibration fit on the validation split
(Table~\ref{tab:valuecal}), losses shrink from $-11.5\%$ to $-7.1\%$
at EV $>0$ but remain strongly negative. The model's apparent value
is concentrated exactly where it is most confident and most wrong:
overconfidence, not information. Consistent with this, betting the
de-vigged opening odds against closing odds (and the reverse)
produces no positive ROI, confirming that opening-to-closing drift
contains no exploitable information in these data. This
value-betting result is specific to \emph{closing} odds: the market's
final line is the hardest benchmark, and the negative finding does
not speak to value-betting strategies evaluated against opening odds
or other pre-closing benchmarks. A fractional Kelly
policy amplifies the losses (a starting bankroll of 100 units falls
to 1.6): stake sizing cannot create an edge that the probabilities do
not have.

% ================= 4.8 error analysis =================
\subsection{Error Analysis: Where the Model Fails}
\begin{table*}[tbp]
\centering
\small
\caption{Error analysis on the test set. Rows 1--3: accuracy of XGBoost
vs de-vigged closing odds on subsets where the two agree or disagree.
Rows 4--6: error rate among predictions with model confidence at least
0.50/0.60/0.70 (N = model/market). Rows 7--10: error rate and share of
strong favourites (min opening odds $\leq$ 1.6) for the extreme
quintiles of odds movement and of market dispersion ((Max-Avg)/Avg
closing odds). Rows 11--12: mean model confidence for draws that were
correctly vs incorrectly classified.}
\label{tab:error}
\resizebox{\textwidth}{!}{
\begin{tabular}{lccc}
\toprule
Group & N & Acc/Err. & Fav. share \\
\midrule
All matches & 1104 & 54.2\% & 54.2\% \\
Model = market & 1034 & 55.6\% & 55.6\% \\
Model $\neq$ market & 70 & 32.9\% & 32.9\% \\
conf $\geq$ 0.50 & 528/524 & 33.9\% & 34.9\% \\
conf $\geq$ 0.60 & 296/259 & 27.0\% & 24.3\% \\
conf $\geq$ 0.70 & 139/100 & 22.3\% & 16.0\% \\
Odds-move Q0 (small) & 229 & 50.7\% & 11\% \\
Odds-move Q4 (large) & 220 & 29.5\% & 66\% \\
Dispersion Q0 (low) & 221 & 52.9\% & 6\% \\
Dispersion Q4 (high) & 221 & 35.7\% & 71\% \\
Draws, correctly predicted & 9 & 0.366 & --- \\
Draws, misclassified & 272 & 0.498 & --- \\
\bottomrule
\end{tabular}
}
\end{table*}


Table~\ref{tab:error} dissects the model's errors on the test set.
First, the model and the market disagree on only 70 of 1,104 matches
(6.3\%), and on those matches the two are equally accurate (32.9\%
each); on the 1,034 agreeing matches they are again equal (55.6\%).
The model therefore contributes no incremental information beyond the
closing line: it neither corrects the market where they disagree nor
adds accuracy on the easy majority where they agree. The low accuracy
on disagreeing matches (32.9\%) further indicates that these are
intrinsically difficult matches, not mispriced ones.

Second, the apparent relationships between error rate and market
dynamics are confounded by market strength. The error rate is lowest
in the largest odds-movement quintile (29.5\%) and the largest
dispersion quintile (35.7\%), but those quintiles are dominated by
strong favourites (66\% and 71\% of matches with minimum opening odds
$\leq$ 1.6, versus 11\% and 6\% in the smallest quintiles), where the
favourite wins most often. Odds movement and bookmaker dispersion are
therefore not independent difficulty signals. The uncertainty index
of Table~\ref{tab:uixgb} is, by contrast, a monotone stratifier
(67.9\% vs 44.4\% accuracy across tiers) that is not reducible to
market strength.

Third, the model is overconfident exactly where it is wrong. The
error rate among predictions with confidence $\geq$ 0.7 is 22.3\%
for the model versus 16.0\% for the market at the same threshold.
For draws---which the model almost never predicts
(Table~\ref{tab:drawclass})---the 272 misclassified draws carry a
mean confidence of 0.498, higher than the 0.366 of the 9 correctly
predicted draws. Confidence is highest precisely when the model is
wrong, consistent with the value-betting result of
Table~\ref{tab:value}: apparent edge is overconfidence, not
information.

% ================= 4.9 SCS/UI stratification =================
\subsection{Uncertainty-Driven Risk Control}
\begin{table}[htbp]
\centering
\small
\caption{Scenario-complexity tiers (XGBoost, test set, thresholds 1.4/1.8). Higher SCS indicates harder matches.}
\label{tab:scs}
\resizebox{\linewidth}{!}{
\begin{tabular}{lccc}
\toprule
SCS tier & N & Acc & ROI \\
\midrule
low(SCS$\leq$1.4) & 246 & 58.9\% & 5.1\% \\
mid(1.4$<$SCS$\leq$1.8) & 356 & 57.3\% & 7.0\% \\
high(SCS$>$1.8) & 502 & 49.6\% & -6.3\% \\
\bottomrule
\end{tabular}
}
\end{table}

\begin{table}[htbp]
\centering
\small
\caption{Uncertainty-index tiers (XGBoost, test set, thresholds 0.30/0.45). High-uncertainty matches are excluded by the no-bet policy; their predictive accuracy is reported separately.}
\label{tab:uixgb}
\resizebox{\linewidth}{!}{
\begin{tabular}{lcccc}
\toprule
UI tier & N & Acc & ROI & MDD \\
\midrule
low & 287 & 67.9\% & 4.1\% & -4.9\% \\
medium & 673 & 50.4\% & 1.9\% & -2.4\% \\
high(no-bet) & 144 & 44.4\% & -13.0\% & -17.6\% \\
\bottomrule
\end{tabular}
}
\end{table}

\begin{table}[htbp]
\centering
\small
\caption{Uncertainty-index tiers for the LLM (test set). The same stratification pattern replicates across model families.}
\label{tab:uillm}
\resizebox{\linewidth}{!}{
\begin{tabular}{lcccc}
\toprule
UI tier & N & Acc & ROI & MDD \\
\midrule
low & 259 & 68.7\% & 3.2\% & -4.4\% \\
medium & 689 & 49.8\% & -0.7\% & -4.5\% \\
high(no-bet) & 156 & 47.4\% & -9.6\% & -13.8\% \\
\bottomrule
\end{tabular}
}
\end{table}


The central positive result of this paper is that uncertainty signals
stratify risk in a monotone and cross-model way. Figure~\ref{fig:ui}
shows the UI distribution and the resulting accuracy by tier.
Table~\ref{tab:scs}
splits the test set by the scenario-complexity score: accuracy falls
from 58.9\% (low complexity) to 49.6\% (high complexity), and ROI
from +7.0\% to $-6.3\%$. Table~\ref{tab:uixgb} shows the same
pattern for the uncertainty index on XGBoost: low-uncertainty matches
reach 67.9\% accuracy and +4.1\% ROI, while the high-uncertainty
group---excluded by the no-bet policy---would have lost 13.0\% ROI.
Table~\ref{tab:uillm} replicates the stratification for the LLM
(68.7\% / 49.8\% / 47.4\%), demonstrating that the signal is a
property of the task and the features, not of a particular model.

\paragraph{Comparison with standard uncertainty measures.}
\begin{table}[htbp]
\centering
\small
\caption{Standard predictive-uncertainty measures vs the uncertainty index (test set, XGBoost probabilities, equal-stake betting, same protocol as Table~\ref{tab:policy}). Entropy $H(p)$ and the expected-Brier term $1-\sum p^2$ are pure functions of the predicted distribution; they stratify accuracy about as well as $1-p_{max}$ but inherit its financial failure mode (selecting low-odds favourites whose ROI is essentially zero at coverage 0.5). The UI adds market-volatility terms and keeps positive ROI at both coverage levels. Bootstrap 95\% CIs for ROI all cross zero. Correlations on the test set: UI vs $1-p_{max}$ 0.51, UI vs entropy 0.52 (see text).}
\label{tab:uncbaselines}
\resizebox{\linewidth}{!}{
\begin{tabular}{lcccccc}
\toprule
Score & Cov & Acc & ROI & ROI 95\% CI & Sharpe & Avg odds \\
\midrule
1 $-\,p_{max}$ & 0.7 & 71.0\% & 1.0\% & [-6.4,8.1] & 0.015 & 1.45 \\
1 $-\,p_{max}$ & 0.5 & 64.3\% & 0.1\% & [-6.1,6.4] & 0.001 & 1.61 \\
Entropy $H(p)$ & 0.7 & 70.7\% & 0.2\% & [-7.0,7.4] & 0.003 & 1.45 \\
Entropy $H(p)$ & 0.5 & 64.7\% & 0.7\% & [-5.8,7.1] & 0.009 & 1.60 \\
1 $-\,\sum p^2$ & 0.7 & 70.7\% & 0.3\% & [-7.1,7.1] & 0.004 & 1.45 \\
1 $-\,\sum p^2$ & 0.5 & 64.5\% & 0.4\% & [-6.0,6.9] & 0.005 & 1.60 \\
UI & 0.7 & 66.5\% & 3.2\% & [-5.4,12.1] & 0.040 & 1.63 \\
UI & 0.5 & 60.7\% & 2.0\% & [-5.8,9.4] & 0.023 & 1.77 \\
Market conf. & 0.7 & 73.2\% & 3.5\% & [-3.7,10.3] & 0.055 & 1.43 \\
Market conf. & 0.5 & 63.8\% & -2.3\% & [-8.4,4.0] & -0.031 & 1.58 \\
\bottomrule
\end{tabular}
}
\end{table}

A natural question is whether the composite index is needed at all,
or whether a standard predictive-uncertainty measure would serve as
well. Table~\ref{tab:uncbaselines} compares four scores under the
identical no-bet protocol: $1 - p_{\max}$, entropy $H(p)$, the
expected-Brier term $1 - \sum p^2$, and the UI. On the test set the
three model-only scores are near-collinear ($r = 0.97$ between
$1 - p_{\max}$ and entropy, $r = 1.00$ between $1 - \sum p^2$ and
entropy): they are the same information in different packaging, and
all three stratify accuracy slightly more sharply than the UI
(low-vs-high tier gap of about 29 percentage points versus 17). The
financial stratification is where they diverge. Entropy and
$1 - \sum p^2$ inherit the failure mode of confidence-based
filtering: at coverage 0.5 they retain low-odds favourites (average
selected odds 1.60--1.61 versus the UI's 1.77) and their
ROI is essentially zero ($+0.4$ to $+0.7\%$), whereas the UI, whose
market-volatility terms break ties among equally confident matches,
keeps positive ROI at both coverage levels ($+3.2\%$ at coverage
0.7, $+2.0\%$ at 0.5). The UI correlates only moderately with the
model-only scores ($r \approx 0.5$): it is not a repackaging of
confidence, and its advantage appears exactly where the paper's
thesis predicts---in decision quality, not in accuracy
stratification. As in Table~\ref{tab:policy}, all ROI bootstrap
intervals cross zero, so the contribution is the stratification
structure rather than demonstrated profit.

\paragraph{Weight and threshold sensitivity.}
The UI weights $(0.4, 0.3, 0.15, 0.15)$ are a fixed preset and the
tier thresholds $(0.30, 0.45)$ are tuned on the validation split. To
rule out overfitting to these choices, the stratification
is recomputed on the validation set under six weight configurations (varying each
component by $\pm 0.1$ while keeping the sum at 1). In every
configuration the low-uncertainty tier retains 60--70\% accuracy and
outperforms the medium tier by 10+ points; only the no-bet coverage
varies (69--551 of 1,752 validation matches) as the relative emphasis
shifts between model confidence and market volatility. The conclusion
that uncertainty tiers stratify accuracy is therefore insensitive to
the specific weights. Likewise, the stop-loss parameters (five
consecutive losses or 10\% drawdown) are not consequential: across
seven tested configurations (from 3 losses/5\% drawdown to 10
losses/15\% drawdown) the UI-tier strategy returns identical ROI
(0.37\%), Sharpe (0.026), and drawdown ($-0.5\%$) on the test set,
because the strategy's equity curve rarely approaches the trigger
thresholds in this season.

% ================= 4.10 strategy comparison =================
\subsection{Betting-Strategy Comparison}
\begin{table*}[tbp]
\centering
\small
\caption{Betting-strategy comparison (XGBoost probabilities, test set, bootstrap 95\% CI for ROI). No-bet matches are excluded from ROI but their coverage is reported.}
\label{tab:strategy}
\resizebox{\textwidth}{!}{
\begin{tabular}{lccccc}
\toprule
Strategy & ROI & Sharpe & MDD & N bets & ROI 95\% CI \\
\midrule
Naive (all bets) & 0.5\% & 0.005 & -2.7\% & 1104 & [-5.3,6.7] \\
Conf. threshold 0.4 & 1.5\% & 0.017 & -2.5\% & 894 & [-4.3,7.6] \\
UI tiers + no-bet & 1.8\% & 0.026 & -2.5\% & 960 & [-2.6,6.2] \\
UI tiers + stop-loss & 0.4\% & 0.026 & -0.5\% & 960 & [-0.5,1.2] \\
\bottomrule
\end{tabular}
}
\end{table*}


Table~\ref{tab:strategy} compares betting policies on the same
XGBoost probabilities, and Figure~\ref{fig:roi} plots the
corresponding cumulative return curves. The naive strategy
(all bets, equal stake)
returns $+0.5\%$; a confidence threshold of 0.4 raises this to
$+1.5\%$; and the UI-tier policy---full stake at low UI, scaled stake
at medium UI, no-bet at high UI---achieves $+1.8\%$ ROI with a
Sharpe ratio of 0.026, while excluding 144 of 1,104 matches. Adding
the stop-loss rule (five consecutive losses or 10\% drawdown from
peak) cuts the maximum drawdown from $-2.7\%$ to $-0.5\%$ at a small
ROI cost. These gains are modest in absolute terms, as they must be
in an efficient market, but they are consistent with the monotone
stratification above and are the difference between a risk-blind
forecaster and a risk-aware decision policy.

% ================= 4.11 LLM =================
\subsection{LLM Analysis}
\label{sec:llm_analysis}
\begin{table}[htbp]
\centering
\small
\caption{Temperature sensitivity (same 120 matches). Multi-sample consistency remains near 1.0 at both temperatures and does not correlate with correctness.}
\label{tab:temp}
\resizebox{\linewidth}{!}{
\begin{tabular}{lccccc}
\toprule
Config & Acc & LogLoss & ECE & Cons. & r(cons,corr) \\
\midrule
DeepSeek t=0.3 & 54.2\% & 0.927 & 0.097 & 0.995 & -0.068 \\
DeepSeek t=0.7 & 55.0\% & 0.927 & 0.089 & 0.990 & 0.037 \\
\bottomrule
\end{tabular}
}
\end{table}

\begin{table}[htbp]
\centering
\small
\caption{Open-weight vs closed API LLM on the same 200 matches. This comparison is robustness evidence only: the sample is small and confidence intervals overlap, so no superiority claim is made for either model.}
\label{tab:open}
\resizebox{\linewidth}{!}{
\begin{tabular}{lccccc}
\toprule
Model & Acc & LogLoss & ECE & ROI & Win rate \\
\midrule
DeepSeek (API) & 56.5\% & 0.934 & 0.065 & 2.0\% & 56.5\% \\
Qwen2.5-Coder-7B (local) & 58.5\% & 0.927 & 0.056 & 7.4\% & 58.5\% \\
\bottomrule
\end{tabular}
}
\end{table}

\begin{table}[htbp]
\centering
\small
\caption{LLM input ablation on the same first 120 test-set matches. Removing all market information degrades accuracy and log loss substantially (49.2\%, 1.028); adding team statistics and referee context to the market signal changes nothing (54.2--55.0\%, log loss 0.922--0.927). The LLM's market-level calibration is therefore driven by the market signal in the prompt, not by independent reasoning over non-market features---consistent with the feature-importance result.}
\label{tab:llmablation}
\resizebox{\linewidth}{!}{
\begin{tabular}{lcccc}
\toprule
Input & Acc & LogLoss & ECE & ROI \\
\midrule
Market only & 55.0\% & 0.922 & 0.087 & -3.5\% \\
Team statistics only & 49.2\% & 1.028 & 0.072 & -7.2\% \\
Market + statistics & 54.2\% & 0.927 & 0.084 & -5.5\% \\
Full card (incl. referee) & 54.2\% & 0.927 & 0.097 & -5.5\% \\
\bottomrule
\end{tabular}
}
\end{table}


The LLM row of Table~\ref{tab:main} summarizes the full test period:
1104 of 1104 matches parsed successfully at a total cost of \$0.55
(3,312 API calls), with accuracy and calibration statistically equal
to the market. Table~\ref{tab:temp} reports the
 temperature
sensitivity: raising temperature from 0.3 to 0.7 changes nothing
material, and multi-sample consistency remains near 1.0 at both
temperatures with no correlation to correctness ($r = -0.07$ and
$+0.04$). The multi-sample consistency mechanism proposed in earlier
work is therefore not informative in this setting; the uncertainty
index derives its power from predictive confidence and market
volatility. Table~\ref{tab:open} compares models on the same first
120 matches of the test set: an open-weight 7B model
(Qwen2.5-Coder-7B, quantized, running on a laptop GPU), the
chat-optimized API model, and a reasoning model (DeepSeek-Reasoner,
run with a single sample because reasoning models are roughly an
order of magnitude slower). The open-weight model matches or slightly
exceeds both API models in accuracy and calibration at zero marginal
cost; the reasoning model improves log loss over chat
(0.906 vs 0.927) but not accuracy, at a cost of \$0.60 for 120
matches---more than the \$0.55 of the entire 1,104-match chat
run. This comparison is \emph{robustness evidence only}: the sample
is small (120 matches), all confidence intervals overlap, and
no claim of superiority is made for any model. Its purpose is to
show that the paper's qualitative conclusions do not depend on a
specific proprietary API or on the choice of reasoning vs
non-reasoning model.

A remaining question is whether the LLM contributes any reasoning
beyond transcribing the market signal it is given. Table~\ref{tab:llmablation}
answers it directly with an input ablation on the same 120 matches:
removing market information degrades the LLM substantially
(49.2\% accuracy, 1.028 log loss), while adding team statistics and
referee context to the market signal changes nothing (54.2--55.0\%,
log loss 0.922--0.927). The LLM's market-level calibration is
therefore driven by the market signal in the prompt, not by
independent reasoning over non-market features---a result that
reinforces the feature-importance finding (Table~\ref{tab:importance})
and keeps the LLM contribution honest: \emph{calibration without
training}, not incremental information. A no-LLM reference point
makes this precise: a baseline that simply outputs the de-vigged
market probabilities achieves 0.55 accuracy and 0.923 log loss on
the same 120 matches (Table~\ref{tab:llmablation}, final row),
statistically indistinguishable from the LLM with the market signal
(0.55, 0.922). Across the full test set, the LLM's output deviates
from the de-vigged market probabilities by a mean absolute error of
0.004 per class (per-class correlations $\geq 0.996$). The LLM is
therefore best characterized as a calibrated transcription interface
for the market signal in its prompt, not as an independent
forecaster.

Finally, a mechanistic explanation for the null
consistency result is offered. In open-domain forecasting settings where
multi-sample consistency is informative, the LLM is typically given
retrieval evidence or heterogeneous free-form context, so different
samples draw on different information. Here the model receives a
\emph{fixed, complete, structured feature card} with no retrieval: at
any temperature the input is identical across samples, the reasoning
path is short, and the output probabilities collapse to a narrow
range (consistency $\geq 0.95$ for 99\% of matches at both
temperatures). Sampling consistency is thus informative only when the
input itself carries sampling variance; in a deterministic structured
task it degrades to a constant and contributes nothing to the
uncertainty index.
